% GENERATED FILE. Edit canonical lessons, then run npm run generate:books.
% canonical-source-hash: fnv1a64:4d9590dc3464cfaf
% canonical-lessons: RU-C481-lopatka, RU-C481-pozvonochnik, RU-C481-bedro, RU-C481-pyatka, RU-C481-ladon

\chapter{From Shoulder Blade to Heel}
\label{ch:ru-from-shoulder-blade-to-heel}
\hlchaptermodality{481}

\begin{chapteropening}
\textbf{By the end of this chapter:} \emph{I can name more parts of the back, the leg and the hand.}

Complete the last lesson of chapter 481: I can name more parts of the back, the leg and the hand.
\end{chapteropening}

\section[\texorpdfstring{\ru{лопатка} (lopátka)}{лопатка (lopátka)}]{\ru{лопатка} (lopátka) — a shoulder blade; a small spade}

\label{lesson:RU-C481-lopatka}



\begin{quote}
Before the new one: say the Russian for to comfort, then the Russian for consolation.
\end{quote}

\subsection*{You'll want to know: \ru{лопатка}}
\textbf{\ru{лопатка}} (\emph{lopátka}) — \textquotedblleft{}a shoulder blade; a small spade\textquotedblright{}.

\begin{quote}
\ru{У} \ru{меня} \ru{болит} \ru{спина} \ru{под} \ru{левой} \ru{лопаткой}.
\end{quote}

\begin{quote}
\emph{U menyá bolít spiná pod lévoy lopátkoy.}
\end{quote}

\begin{quote}
My back hurts under my left shoulder blade.
\end{quote}

Feminine. Also a small spade or a trowel: a diminutive of \emph{\ru{лопата}} — a spade.

\begin{cousinweb}
\emph{\ru{лопатка}} is a diminutive of \emph{\ru{лопата}} 'spade, shovel', which is inherited from Proto-Slavic \emph{lopata} (EtymDB); the flat bone was named after the flat blade.
\end{cousinweb}

\begin{grammarlens}[title={What you've built}]
One more word for telling, explaining and describing at B1.
\end{grammarlens}

\subsection*{Your turn}
\begin{itemize}
\raggedright
  \item \emph{Say it:} \emph{lopátka}
  \item \emph{Say it:} \emph{lopátka}, once more
  \item \emph{Say it:} a sentence of your own with \emph{lopátka}
  \item \emph{Recall:} say the Russian for to comfort, then the Russian for consolation, then say \emph{lopátka} again
\end{itemize}

\begin{tcolorbox}[breakable,colback=teal!4,colframe=teal!35!black,title={Before you move on}]
What does \ru{лопатка} mean? (\textquotedblleft{}A shoulder blade; a small spade\textquotedblright{}.) Say it once more.
\end{tcolorbox}
\section[\texorpdfstring{\ru{позвоночник} (pozvonóchnik)}{позвоночник (pozvonóchnik)}]{\ru{позвоночник} (pozvonóchnik) — the spine}

\label{lesson:RU-C481-pozvonochnik}



\begin{quote}
Before the new one: say the Russian for consolation, then the Russian for a shoulder blade.
\end{quote}

\subsection*{You'll want to know: \ru{позвоночник}}
\textbf{\ru{позвоночник}} (\emph{pozvonóchnik}) — \textquotedblleft{}the spine\textquotedblright{}.

\begin{quote}
\ru{Врач} \ru{сказал}, \ru{что} \ru{у} \ru{меня} \ru{здоровый} \ru{позвоночник}.
\end{quote}

\begin{quote}
\emph{Vrach skazál, chto u menyá zdoróvyy pozvonóchnik.}
\end{quote}

\begin{quote}
The doctor said that my spine is healthy.
\end{quote}

Masculine. Built on \emph{\ru{позвонок}} — a vertebra.

\begin{cousinweb}
\emph{\ru{позвоночник}} is built on \emph{\ru{позвонок}} 'vertebra' with \emph{-\ru{ник}}; EtymDB has no entry, and the origin of \emph{\ru{позвонок}} is not settled.
\end{cousinweb}

\begin{grammarlens}[title={What you've built}]
One more word for telling, explaining and describing at B1.
\end{grammarlens}

\subsection*{Your turn}
\begin{itemize}
\raggedright
  \item \emph{Say it:} \emph{pozvonóchnik}
  \item \emph{Say it:} \emph{pozvonóchnik}, once more
  \item \emph{Say it:} a sentence of your own with \emph{pozvonóchnik}
  \item \emph{Recall:} say the Russian for consolation, then the Russian for a shoulder blade, then say \emph{pozvonóchnik} again
\end{itemize}

\begin{tcolorbox}[breakable,colback=teal!4,colframe=teal!35!black,title={Before you move on}]
What does \ru{позвоночник} mean? (\textquotedblleft{}The spine\textquotedblright{}.) Say it once more.
\end{tcolorbox}
\section[\texorpdfstring{\ru{бедро} (bedró)}{бедро (bedró)}]{\ru{бедро} (bedró) — a thigh; the hip}

\label{lesson:RU-C481-bedro}



\begin{quote}
Before the new one: say the Russian for a shoulder blade, then the Russian for the spine.
\end{quote}

\subsection*{You'll want to know: \ru{бедро}}
\textbf{\ru{бедро}} (\emph{bedró}) — \textquotedblleft{}a thigh; the hip\textquotedblright{}.

\begin{quote}
\ru{Он} \ru{упал} \ru{и} \ru{ударился} \ru{бедром}.
\end{quote}

\begin{quote}
\emph{On upál i udárilsya bedróm.}
\end{quote}

\begin{quote}
He fell and hit his hip.
\end{quote}

Neuter; plural \ru{бёдра}.

\begin{cousinweb}
\emph{\ru{бедро}} is inherited from Proto-Slavic \emph{bedro}, which EtymDB traces back to PIE \emph{b\textsuperscript{h}ed\textsuperscript{h}ros.}
\end{cousinweb}

\begin{grammarlens}[title={What you've built}]
One more word for telling, explaining and describing at B1.
\end{grammarlens}

\subsection*{Your turn}
\begin{itemize}
\raggedright
  \item \emph{Say it:} \emph{bedró}
  \item \emph{Say it:} \emph{bedró}, once more
  \item \emph{Say it:} a sentence of your own with \emph{bedró}
  \item \emph{Recall:} say the Russian for a shoulder blade, then the Russian for the spine, then say \emph{bedró} again
\end{itemize}

\begin{tcolorbox}[breakable,colback=teal!4,colframe=teal!35!black,title={Before you move on}]
What does \ru{бедро} mean? (\textquotedblleft{}A thigh; the hip\textquotedblright{}.) Say it once more.
\end{tcolorbox}
\section[\texorpdfstring{\ru{пятка} (pyátka)}{пятка (pyátka)}]{\ru{пятка} (pyátka) — a heel}

\label{lesson:RU-C481-pyatka}



\begin{quote}
Before the new one: say the Russian for the spine, then the Russian for a thigh.
\end{quote}

\subsection*{You'll want to know: \ru{пятка}}
\textbf{\ru{пятка}} (\emph{pyátka}) — \textquotedblleft{}a heel\textquotedblright{}.

\begin{quote}
\ru{Новые} \ru{туфли} \ru{натёрли} \ru{мне} \ru{пятку}.
\end{quote}

\begin{quote}
\emph{Nóvyye túfli natyórli mne pyátku.}
\end{quote}

\begin{quote}
My new shoes have rubbed my heel.
\end{quote}

Feminine. A diminutive in form of the old word \emph{\ru{пята}}, now the everyday word.

\begin{cousinweb}
\emph{\ru{пятка}} is a diminutive in form of \emph{\ru{пята}} 'heel', which is inherited from Proto-Slavic \emph{pęta} 'heel' (EtymDB).
\end{cousinweb}

\begin{grammarlens}[title={What you've built}]
One more word for telling, explaining and describing at B1.
\end{grammarlens}

\subsection*{Your turn}
\begin{itemize}
\raggedright
  \item \emph{Say it:} \emph{pyátka}
  \item \emph{Say it:} \emph{pyátka}, once more
  \item \emph{Say it:} a sentence of your own with \emph{pyátka}
  \item \emph{Recall:} say the Russian for the spine, then the Russian for a thigh, then say \emph{pyátka} again
\end{itemize}

\begin{tcolorbox}[breakable,colback=teal!4,colframe=teal!35!black,title={Before you move on}]
What does \ru{пятка} mean? (\textquotedblleft{}A heel\textquotedblright{}.) Say it once more.
\end{tcolorbox}
\section[\texorpdfstring{\ru{ладонь} (ladón)}{ладонь (ladón)}]{\ru{ладонь} (ladón) — a palm (of the hand)}

\label{lesson:RU-C481-ladon}



\begin{quote}
Before the new one: say the Russian for a thigh, then the Russian for a heel.
\end{quote}

\subsection*{You'll want to know: \ru{ладонь}}
\textbf{\ru{ладонь}} (\emph{ladón}) — \textquotedblleft{}a palm (of the hand)\textquotedblright{}.

\begin{quote}
\ru{Ребёнок} \ru{держал} \ru{бабочку} \ru{на} \ru{ладони}.
\end{quote}

\begin{quote}
\emph{Rebyónok derzhál bábochku na ladóni.}
\end{quote}

\begin{quote}
The child held a butterfly on his palm.
\end{quote}

Feminine.

\begin{cousinweb}
EtymDB lists \emph{\ru{ладонь}} without a source; it is usually traced to older \emph{\ru{долонь}}, from Proto-Slavic \emph{doln\ru{ь}}. Church Slavonic \emph{\ru{длань}} is the same word with \emph{-\ru{ла}-}: EtymDB traces \emph{\ru{длань}} through Church Slavonic to Proto-Slavic \emph{doln\ru{ь}} 'palm'.
\end{cousinweb}

\begin{grammarlens}[title={What you've built}]
One more word for telling, explaining and describing at B1.
\end{grammarlens}

\subsection*{Your turn}
\begin{itemize}
\raggedright
  \item \emph{Say it:} \emph{ladón}
  \item \emph{Say it:} \emph{ladón}, once more
  \item \emph{Say it:} a sentence of your own with \emph{ladón}
  \item \emph{Recall:} say the Russian for a thigh, then the Russian for a heel, then say \emph{ladón} again
\end{itemize}

\begin{tcolorbox}[breakable,colback=teal!4,colframe=teal!35!black,title={Before you move on}]
What does \ru{ладонь} mean? (\textquotedblleft{}A palm (of the hand)\textquotedblright{}.) Say it once more.
\end{tcolorbox}
